
Suponga que una vía de ferrocarril se construyó con rieles de cierta longitud \textit{L}, dejando entre ellos juntas de dilatación de 1cm de amplitud.
\begin{enumerate}
    \item Si la vía se construyera con rieles de mayor longitud que \textit{L}, ¿las juntas de dilatación deberán tener una amplitud mayor, menor o igual a 1 cm? Explique.
    \item ¿Por qué, si se produce un incendio sobre una vía férrea, los rieles se deforman a pesar de la existencia de las juntas de dilatación?
\end{enumerate}